\documentclass[11pt]{article}
\usepackage[margin=1in]{geometry}
\usepackage{amsmath,amssymb}
\usepackage{graphicx}
\usepackage{booktabs}
\usepackage{siunitx}
\usepackage{xcolor}
\usepackage[colorlinks=true,linkcolor=blue,urlcolor=blue]{hyperref}

\newcommand{\shift}{\ensuremath{\Delta N}}
\newcommand{\diff}{\ensuremath{\delta}}
\newcommand{\lm}{\ensuremath{\lambda}}
\newcommand{\mumu}{\ensuremath{\mu}}

\title{\bf Selection Rules for the U(3)$\times$U(4) Tensor Expansion\\
of Two-Body Operators}
\author{SAM --- Symmetry-Adapted Many-body framework}
\date{\today}

\begin{document}
\maketitle

\begin{abstract}
When a two-body operator is expanded into $\mathrm{U}(3)\times\mathrm{U}(4)$
(space $\times$ Wigner spin--isospin) tensors, roughly half of the tensors that
pass the current kinematic filters are found to have vanishing matrix elements.
This note collects the selection rules that explain that structure, derives them
from harmonic-oscillator $\mathrm{U}(3)$ and Wigner $\mathrm{SU}(4)$ group
theory, and records their empirical validation across the chiral-EFT interaction
family. The central results are: for a \emph{central} (rotational-scalar)
two-body force the terminal $\mathrm{SU}(3)$ irrep obeys $\lm-\mumu=\shift$
exactly and $\lm+\mumu$ even (translational invariance and parity), and its
$\mathrm{SU}(4)$ content is
restricted to the four self-conjugate irreps
$\{(0,0,0),(0,2,0),(1,0,1),(2,0,2)\}$. Both are
consequences of the force preserving the two-particle relative/exchange
structure; a normal-ordered three-body force breaks both through the contracted
core nucleon, in a controlled way that \emph{tags} the induced content.
\end{abstract}

\section{Setup and notation}
Single-particle creation operators in oscillator shell $n$ transform as the
symmetric $\mathrm{SU}(3)$ irrep $(n,0)$ and the $\mathrm{SU}(4)$ fundamental
$\mathbf 4=(1,0,0)$; annihilation operators transform as the conjugates $(0,n)$
and $\bar{\mathbf 4}=(0,0,1)$. A two-body operator
$a^\dagger_{n_a}a^\dagger_{n_b}\,\tilde a_{n_c}\tilde a_{n_d}$, coupled to good
$\mathrm{U}(3)\times\mathrm{U}(4)$, terminates in an irrep whose labels we write
as
\[
\text{U(3): } [n_{13},n_{23},n_{33}],\quad
(\lm,\mumu)=(n_{13}-n_{23},\,n_{23}-n_{33}); \qquad
\text{U(4): } [n_{14},n_{24},n_{34},n_{44}],\quad
(P_1,P_2,P_3).
\]
The net oscillator-quanta change carried by the operator is
\[
\shift \;=\; \text{shift\_quanta} \;=\; (n_a+n_b)-(n_c+n_d)
      \;=\;\sum_{\text{shells}} n\,(n_{\text{ads}}-n_{\text{ats}}).
\]
We define the U(3) \emph{defect}
\[
\boxed{\ \diff \;\equiv\; (\lm-\mumu)-\shift\ }.
\]

\section{U(3) sector}

\subsection{Empirical rules}
For every \emph{genuine, central} two-body force examined,
\[
\diff = 0, \qquad \lm+\mumu \ \text{even}
\qquad\text{(exactly, for all matrix-element--nonzero tensors).}
\]
The first is the central/relative-motion rule ($h_2=0$); the second is parity
conservation. Universally, for \emph{any} operator (scalar or not, two- or
three-body),
\[
\diff \equiv 0 \pmod 3 .
\]

\subsection{The governing identity}
The operator is a $\mathrm{U}(3)$ tensor with a top Gelfand--Tsetlin row
$[h_1,h_2,h_3]$. Its $\mathrm{U}(1)$ charge (number of quanta) is
$h_1+h_2+h_3=\shift$, and $(\lm,\mumu)=(h_1-h_2,\,h_2-h_3)$. Hence
\[
\lm-\mumu=(h_1-h_2)-(h_2-h_3)=(h_1+h_2+h_3)-3h_2=\shift-3h_2,
\]
so that
\[
\boxed{\ \diff = -3\,h_2\ }
\qquad (h_2 = \text{middle entry of the operator's top Gelfand row}).
\]
Two corollaries follow immediately:
\begin{itemize}
\item $\diff=-3h_2$ is \emph{always} a multiple of $3$ --- this is the
$\mathrm{SU}(3)$ triality law $\diff\equiv0\pmod3$; a defect not divisible by
$3$ is group-theoretically impossible.
\item $\diff=0 \iff h_2=0$: the operator lives in a two-rowed, maximally
\emph{stretched} irrep $[h_1,0,h_3]$.
\end{itemize}

\subsection{Why translational invariance forces $h_2=0$}
The single property at work is \emph{translational invariance}, not centrality.
A translationally invariant two-body operator depends only on the relative
coordinate $\mathbf r_1-\mathbf r_2$; the center of mass is a spectator:
\[
V=(\text{relative-motion operator})\otimes(\text{CM identity}).
\]
The relative coordinate is a \emph{single} 3D oscillator, so any operator on it
is a polynomial in $a^\dagger_{\rm rel},\tilde a_{\rm rel}$: a term
$a^{\dagger\,p}_{\rm rel}\tilde a^{\,q}_{\rm rel}$ couples
$(p,0)\otimes(0,q)=\bigoplus_k (p-k,\,q-k)$, and every piece has
$\lm-\mumu=p-q$ --- as a $\mathrm{U}(3)$ irrep of charge $p-q$ it is two-rowed,
$h_2=0$. (Rotational rank is irrelevant: a translationally invariant
\emph{tensor} operator, e.g.\ the relative quadrupole, is equally two-rowed.)
The CM identity is a $\mathrm{U}(3)$ scalar $(0,0)$, so coupling to it keeps
$h_2=0$, and since the CM carries no quanta $\shift_{\rm rel}=\shift$. Therefore
\[
\lm-\mumu=\shift .\qquad\blacksquare
\]
A central potential $V(|\mathbf r_1-\mathbf r_2|)$ is one such operator, but the
result holds for any translationally invariant two-body operator. Conversely a
tensor with $h_2\neq0$ (i.e.\ $\diff\neq0$) requires the CM coordinate to
participate --- the operator is not of the form $V_{\rm rel}\otimes I_{\rm CM}$
--- which is exactly the population of ``accepted-then-empty'' tensors.

\subsection{Translational invariance and parity: $\lm+\mumu$ even}
The terminal label is $n_{13}=\lm+\mumu$, and the filter requires
\[
\boxed{\ \lm+\mumu \ \text{even}\ } \qquad\text{(equivalently } \lm,\mumu \text{ share parity).}
\]
From the top Gelfand row, $\lm+\mumu=h_1-h_3$ and $\shift=h_1+h_2+h_3$, so
\[
\lm+\mumu \equiv \shift + h_2 \pmod 2,
\qquad\text{hence}\qquad
\lm+\mumu\ \text{even}\iff \shift\ \text{even}\ \wedge\ h_2\ \text{even}.
\]
The two conditions are two distinct symmetries:
\begin{itemize}
\item $\shift$ even is \textbf{parity}. The operator carries oscillator parity
$(-1)^{\shift}$ (it changes the quanta by $\shift$); a parity-even scalar force
has $\shift$ even.
\item $h_2$ even is \textbf{translational invariance}. A translationally
invariant two-body force acts only on the relative coordinate with the CM as a
spectator, which forces $h_2=0$ (Sec.~3.3, the same origin as $\diff=-3h_2=0$);
in particular $h_2$ is even. So $\lm+\mumu$ even is the mod-$2$ shadow of
$\lm-\mumu=\shift$, and for a genuine two-body scalar the latter already implies
it.
\end{itemize}
A nuclear force \emph{always} conserves parity, so $\shift$ is even for both
genuine two-body potentials and normal-ordered $2\mathrm B{+}3\mathrm B$
operators; parity is therefore \emph{not} what distinguishes them. The
discriminator is $h_2$: translational invariance pins $h_2=0$ (even), whereas a
normal-ordered three-body force \textbf{breaks translational invariance} --- the
contracted core nucleon is not a CM spectator --- and admits $h_2\in\{0,\pm1,
\pm2,\pm3\}$ (from $\diff=-3h_2\in\{0,\pm3,\pm6,\pm9\}$). Odd $h_2$ then gives
$\lm+\mumu$ \emph{odd} with parity fully intact. Empirically HO\_Be7/Li7 keep
many such tensors ($4688$ of $13164$ at Nmax6; $115752$ of $280354$ at Nmax10),
so $\lm+\mumu$ even is a two-body-structure rule gated to recognized scalar
potentials and must \emph{not} be imposed on NO-3B.

The rule is also independent of triality ($\diff\equiv0\pmod3$ permits
$\lm-\mumu=3$, i.e.\ $\lm+\mumu$ odd), and it is what makes the only realized
parity classes $(\lm,\mumu)\bmod2\in\{(0,0),(1,1)\}$ --- the backdrop against
which the $\mathrm{SU}(4)$-singlet rule (both even) sharpens to pick out the
space-symmetric sector.

\subsection{Non-scalar and three-body deviations}
\begin{itemize}
\item \textbf{Non-scalar operators} (e.g.\ Gamow--Teller, $JJ=2$): the
spin/orbital tensor structure couples CM and relative motion, so $h_2\neq0$
components acquire strength and $\diff=-3h_2\in\{0,\pm3,\pm6,\pm9\}$.
\item \textbf{Normal-ordered three-body forces}: the contracted third nucleon
is a \emph{core} orbital, not a CM spectator, and injects $h_2\neq0$. Even a
$J$-scalar Hamiltonian (e.g.\ HO\_Be7/Li7) then shows
$\diff\in\{\pm3,\pm6,\pm9\}$ with substantial strength. Thus, \emph{for a scalar
operator}, $\diff\neq0$ with nonzero strength \emph{tags} the three-body-induced
content, with no strength evaluation needed.
\end{itemize}

\section{U(4) (Wigner supermultiplet) sector}

\subsection{Empirical rule}
For a central two-body force, the terminal $\mathrm{U}(4)$ irrep is restricted
to the four \emph{self-conjugate} $\mathrm{SU}(4)$ irreps:
\[
\begin{array}{llll}
\text{[0,0,0,0]} & (0,0,0) & \mathbf 1 \\
\text{[2,2,0,0]} & (0,2,0) & \mathbf{20} \\
\text{[2,1,1,0]} & (1,0,1) & \mathbf{15} \\
\text{[4,2,2,0]} & (2,0,2) & \mathbf{84}
\end{array}
\qquad\text{(self-conjugate}\iff P_1=P_3\text{).}
\]

\subsection{Group-theoretic origin}
The two constituent pairs decompose as
\[
\mathbf 4\otimes\mathbf 4=\underbrace{\mathbf{10}}_{(2,0,0)\text{ sym}}
   \oplus\underbrace{\mathbf 6}_{(0,1,0)\text{ antisym}},\qquad
\bar{\mathbf 4}\otimes\bar{\mathbf 4}=\overline{\mathbf{10}}\oplus\mathbf 6 ,
\]
($\mathbf 6$ is self-conjugate). The operator
$a^\dagger a^\dagger\,\tilde a\tilde a$ is (creation pair)$\otimes$(annihilation
pair), i.e.\ four products:
\[
\begin{aligned}
\mathbf{10}\otimes\overline{\mathbf{10}} &= (0,0,0)\oplus(1,0,1)\oplus(2,0,2)
   &&\text{(all self-conjugate)},\\
\mathbf 6\otimes\mathbf 6 &= (0,0,0)\oplus(1,0,1)\oplus(0,2,0)
   &&\text{(all self-conjugate)},\\
\mathbf{10}\otimes\mathbf 6 &= (2,1,0)\oplus(1,0,1)
   &&[\,\mathbf{45}\oplus\mathbf{15}\,],\\
\mathbf 6\otimes\overline{\mathbf{10}} &= (0,1,2)\oplus(1,0,1)
   &&[\,\overline{\mathbf{45}}\oplus\mathbf{15}\,].
\end{aligned}
\]
The union of the two \emph{diagonal} products is exactly the four self-conjugate
irreps above. The two \emph{cross} products contribute, beyond the (already
self-conjugate) $(1,0,1)$, the only two \emph{non}-self-conjugate irreps that a
two-body operator can carry:
\[
(2,1,0)=\mathbf{45}\to[3,1,0,0],\quad
(0,1,2)=\overline{\mathbf{45}}\to[3,3,2,0].
\]
These are exactly the extra $\mathrm{U}(4)$ types seen in the normal-ordered
$2\mathrm{B}{+}3\mathrm{B}$ Hamiltonians (HO\_Be7, HO\_Li7), and are absent from
the pure two-body potentials.

\subsection{Why the diagonal (self-conjugate) part only}
Fermion antisymmetry ties $\mathrm{SU}(4)$ type to spatial exchange symmetry:
\[
\mathbf{10}\ (\text{SU(4) sym})\leftrightarrow\text{spatial antisym},\quad
\mathbf 6\ (\text{SU(4) antisym})\leftrightarrow\text{spatial sym}.
\]
A central $V(|\mathbf r_1-\mathbf r_2|)$ commutes with the exchange operator
$P_{12}$, so it cannot change the spatial exchange symmetry: it connects
sym$\to$sym and antisym$\to$antisym only, i.e.\ $\mathbf 6\to\mathbf 6$ and
$\mathbf{10}\to\mathbf{10}$. That selects the diagonal products
$\mathbf 6\otimes\mathbf 6$ and $\mathbf{10}\otimes\overline{\mathbf{10}}$ ---
the four self-conjugate irreps. The cross terms $\mathbf{10}\otimes\mathbf 6$
and $\mathbf 6\otimes\overline{\mathbf{10}}$ --- which alone can produce the
non-self-conjugate $(2,1,0)$ and $(0,1,2)$ --- would require flipping the
spatial exchange symmetry, which a central force is forbidden to do.

\subsection{Three-body deviations}
The 2-body piece of a normal-ordered 3-body force is
$\sum_{\rm core}\langle(12)c|W|(12)c\rangle$: the third leg is contracted against
a core nucleon, which carries its own spin--isospin and coordinate. This
contraction does not commute with $P_{12}$ the way $V(r_{12})$ does; it provides
the handle to connect $\mathbf{10}\leftrightarrow\mathbf 6$, populating the
non-self-conjugate $\mathrm{U}(4)$ irreps. (This is the spin--isospin face of the
same ``the third nucleon is not a spectator'' that gives $\diff\neq0$ in U(3).)

\section{Unified statement and practical filters}
The spatial rules follow from \emph{translational invariance} (the operator is
$V_{\rm rel}\otimes I_{\rm CM}$, forcing $h_2=0$, i.e.\ $\lm-\mumu=\shift$, and
with parity $\lm+\mumu$ even) and the spin--isospin rules from \emph{fermion
antisymmetry} (self-conjugate U(4), and the per-pair/supermultiplet rules
below). A normal-ordered three-body force breaks translational invariance
through the contracted core nucleon, which is why the spatial two-body-structure
rules (R0, R1) --- and, in the spin--isospin sector, R4 --- must be switched off
for it, while the antisymmetrization rules (R2, R3) survive.

\begin{center}
\small
\setlength{\tabcolsep}{5pt}
\resizebox{\textwidth}{!}{%
\begin{tabular}{@{}lll@{}}
\toprule
Operator class & U(3) rule & U(4) content \\
\midrule
Pure central 2-body & $\lm-\mumu=\shift$ ($\diff=0$), $\lm+\mumu$ even
   & self-conjugate only: $\{(0,0,0),(0,2,0),(1,0,1),(2,0,2)\}$ \\
Normal-ordered 2B+3B & $\diff\equiv0\pmod3$ ($\lm+\mumu$ may be odd)
   & self-conjugate $+\ (2,1,0),(0,1,2)$ \\
Non-scalar (rank $J_0$) & $\diff\equiv0\pmod3,\ |\diff|\le 3f(J_0)$
   & rank-dependent \\
\bottomrule
\end{tabular}}
\end{center}

\noindent As a pre-filter for the scalar-potential workload, $\diff=0$ safely
removes the $\mathrm{U}(3)$-forbidden tensors before their (expensive) matrix
elements are computed; a nonzero-strength tensor with $\diff\neq0$ in a scalar
operator is either numerical noise ($\diff\in6\mathbb{Z}$, tiny strength) or
normal-ordered three-body content ($\diff\in\{\pm3,\pm9,\dots\}$).

\section{The implemented filter rules}
The expansion pre-filter (\texttt{TwoBodyOperator::passes\_selection}) applies a
parity rule (R0, in \texttt{U3selection}) and four further exact,
strength-independent rules (R1--R4), alongside the $\mathrm{U}(4)$
self-conjugate restriction of Sec.~4 (in \texttt{U4selection}). Each was
validated with \emph{zero false negatives} (no matrix-element--nonzero tensor is
ever rejected). R0 is applied upstream of the matrix-element stage on which the
coverage statistics are collected, so the $96.6\%$ figure below is the fraction
of the remaining (parity-even) ``accepted-then-empty'' tensors removed by
R1--R4. Write the terminal irrep as $(\lm,\mumu)$ and $(P_1,P_2,P_3)$, and
recall the per-shell constituent pairs: a same-shell creation pair transforms in
$(n,0)\otimes(n,0)$, an annihilation pair in the conjugate $(0,n)\otimes(0,n)$,
with $\mathrm{SU}(4)$ either the symmetric $\mathbf{10}=(2,0,0)$ /
$\overline{\mathbf{10}}=(0,0,2)$ or the antisymmetric $\mathbf 6=(0,1,0)$.

\begin{description}
\item[R0 --- Translational invariance $+$ parity ($\lm+\mumu$ even).] The
terminal label $n_{13}=\lm+\mumu$ must be even (equivalently $\lm,\mumu$ share
parity). Since $\lm+\mumu\equiv\shift+h_2\pmod2$, this is $\shift$ even (parity)
$\wedge$ $h_2$ even (translational invariance, $h_2=0$); it is the mod-$2$
shadow of R1. Parity holds for NO2B too, so its failure there flags the broken
translational invariance (Sec.~3.4).
\item[R1 --- Translational invariance ($\lm-\mumu=\shift$).] Terminal defect
$\diff=0$, equivalently $h_2=0$ (Sec.~3.3): a translationally invariant two-body
operator is $V_{\rm rel}\otimes I_{\rm CM}$ and carries no CM excitation.
\item[R2 --- Per-pair Pauli antisymmetry.] Every \emph{same-shell} two-particle
(or two-hole) factor must be overall antisymmetric:
\[
\boxed{\ (-1)^{\sigma_{\rm space}}\,(-1)^{\sigma_{\rm SU(4)}}=-1\ },\qquad
(-1)^{\sigma_{\rm space}}=\begin{cases}(-1)^{\mumu_{\rm pair}} & \text{creation }(n,0)^{2}\\[2pt]
(-1)^{\lm_{\rm pair}} & \text{annihilation }(0,n)^{2}\end{cases}
\]
with $(-1)^{\sigma_{\rm SU(4)}}=+1$ for $\mathbf{10}/\overline{\mathbf{10}}$ and
$-1$ for $\mathbf 6$.
\item[R3 --- $\mathrm{SU}(4)$-singlet parity.] If the terminal $\mathrm{SU}(4)$
is the singlet $(0,0,0)$ (spin--isospin scalar channel) then
$\lm$ and $\mumu$ are \emph{both even}. The singlet cannot supply the
spin--isospin exchange that an odd spatial irrep would require.
\item[R4 --- $S{+}T$ parity match.] When a tensor contains both a same-shell
creation pair and annihilation pair, their $\mathrm{SU}(4)$ exchange types must
agree (both $\mathbf 6$, or both $\mathbf{10}/\overline{\mathbf{10}}$). Since
$\mathbf 6$ carries $S{+}T$ odd and $\mathbf{10}/\overline{\mathbf{10}}$ carries
$S{+}T$ even, a central 2-body force conserves the parity of $S+T$.
\end{description}

\paragraph{Two-tier gating.} R2 and R3 are \emph{antisymmetrization /
supermultiplet} rules and hold for \emph{any} $J{=}0$ nuclear interaction,
including the normal-ordered $2\mathrm B{+}3\mathrm B$ operators; they are gated
to all such operators. R0, R1 and R4 are \emph{genuine-2-body-structure} rules:
the contracted core nucleon of a 3-body force breaks all three, so they are
gated to recognized scalar potentials only. Applying them to NO-3B would discard
real content (R0: $115752$, R1: $165929$, R4: $2734$ real losses on HO\_Be7/Li7
at Nmax10), whereas R2/R3 lose nothing there.

\begin{center}
\small
\setlength{\tabcolsep}{6pt}
\begin{tabular}{@{}llcc@{}}
\toprule
Rule & Type & Applies to & Ignores caught \\
\midrule
R0 $\ \lm+\mumu$ even & 2-body structure & potentials & upstream \\
R1 $\ \lm-\mumu=\shift$ & 2-body structure & potentials & 3010 \\
R2 per-pair Pauli & antisymmetrization & potentials $+$ NO2B & 5240 \\
R3 singlet $\Rightarrow\lm,\mumu$ even & antisymmetrization & potentials $+$ NO2B & 1604 \\
R4 $S{+}T$-parity match & 2-body structure & potentials & 114 \\
\midrule
\multicolumn{3}{@{}l}{combined (0 false negatives)} & $9968/10324=96.6\%$ \\
\bottomrule
\end{tabular}
\end{center}

\subsection{The residual: a recoupling boundary}
The remaining $3.4\%$ is not numerical noise --- at Nmax6 the empty tensors are
reproduced \emph{exactly} by their labels (identically for N2LO\_opt and
N2LO\_sat). They are dominated by density-chain structures
$\rho_i\otimes\rho_j=a^\dagger_i a^\dagger_j\tilde a_j\tilde a_i$ built from
$\mathrm{SU}(4)$-adjoint densities. Their vanishing is the same antisymmetry
principle as R2/R4 applied at the \emph{particle--hole (recoupling) level}: the
$\mathrm{SU}(3)$ spatial coupling and the $\mathrm{SU}(4)$ spin--isospin coupling
of the two density factors must carry matching exchange symmetry, or the direct
and exchange contributions cancel. Unlike R1--R4 this is not a one-line label
inequality: it is the structural zero of the $\mathrm{SU}(3)\times\mathrm{SU}(4)$
recoupling (Pandya / $9$-$(\lm\mumu)$) coefficient, which depends jointly on the
factor irreps, the outer multiplicity, and the result. It is left out of the
pre-filter because evaluating that coefficient costs comparably to the matrix
element it would guard.

\section{Numerical validation}
The U(3) rule $\lm-\mumu=\shift$ was verified on the accepted-tensor sets of
the interaction library (Nmax as indicated). ``Real violations'' counts tensors
with $\diff\neq0$ and $\max|\text{strength}|>10^{-6}$.

\begin{center}
\small
\setlength{\tabcolsep}{5pt}
\resizebox{\textwidth}{!}{%
\begin{tabular}{@{}llccl@{}}
\toprule
Operator & $JJ$ & \#tensors & real viol. & note \\
\midrule
N2LO\_opt / sat / srg2.0 (Nmax8) & 0 & 34{,}273 & 0 & scalar 2B \\
N3LO / N4LO chiral, bare \& SRG (Nmax8) & 0 & $\sim$34{,}273 & 0 & scalar 2B \\
N2LO\_opt (Nmax10 / Nmax12) & 0 & 86{,}575 / 256{,}854 & 0 / 0 & noise only ($\sim10^{-8}$) \\
Kinetic, Rrel, CoMNumberOp & 0 & $\le609$ & 0 & scalar \\
\midrule
Gamow--Teller ($N^3$LO EM500) & 2 & 11{,}420 & --- & $\diff\in\{0,\pm3,\pm6,\pm9\}$ (rank) \\
HO\_Be7 / HO\_Li7 (2B $+$ NO-3B) & 0 & 13{,}164 & --- & $\diff\in\{0,\pm3,\pm6,\pm9\}$ (3-body) \\
\bottomrule
\end{tabular}}
\end{center}

\noindent The apparent $\le0.2\%$ ``violations'' at Nmax10/12 are floating-point
noise: their strengths (median $\sim1.6\times10^{-8}$) are $7$--$9$ orders of
magnitude below genuine tensors (median $\sim0.34$), and they still obey
$\diff\equiv0\pmod3$. The two interactions N2LO\_opt and N2LO\_sat give
byte-identical accept/reject statistics, confirming the rule is kinematic (label
dependent), not dynamical (strength dependent). A controlled cross-check:
\texttt{N3LO\_EM500\_srg2.0} and \texttt{HO\_Li7} share $JJ=0,\,TT=2$, yet the
genuine 2-body force has $\diff=0$ only while the NO-3B Hamiltonian reaches
$\diff=\pm3,\pm6,\pm9$ --- isolating the three-body origin.

\end{document}
